\section{Results}
\label{sec:results}
\input{figures/TABLE_T2_matrix.tex}

\begin{figure*}[!htb]
    \centering
    \begin{minipage}[t]{0.48\textwidth}
        \centering
        \includegraphics[width=\linewidth]{figures/fig2_ladder.pdf}
        \Description{Two bar panels. The left panel shows per-iteration time dropping
        from 11.667 milliseconds for the eager reference to 0.619 for the compiled
        dataflow and 0.397 for the CUDA graph, with arrows marking the 19 times, 1.6
        times, and overall 29 times reductions. The right panel shows Aten operations
        per iteration dropping from 1742 to 96, an 18 times reduction.}
        \captionof{figure}{Per-iteration time (left) and PyTorch/Aten operations (right)
        on the HUNL turn subgame.}
        \label{fig:ladder}
    \end{minipage}\hfill
    % 0.465 against fig2's 0.48 so fig6's 3.2in canvas scales by the same 0.80 as
    % fig2's 3.3in one; matching the scale factors keeps the type in both plots at
    % one size, and the remaining 0.055 is the gutter between the two captions.
    \begin{minipage}[t]{0.465\textwidth}
        \centering
        \includegraphics[width=\linewidth]{figures/fig6_scaling.pdf}
        \Description{Log-log plot of milliseconds per CFR-plus iteration against the
        number of States. The measured solid curve rises only 2.9 times while a dash-dot
        reference line of slope one rises 37.9 times; a dashed second axis shows
        achieved bandwidth reaching about 107 gigabytes per second at the largest
        size.}
        \captionof{figure}{HUNL turn scaling, $K=4$ to $24$ hands per player (46k to
        1.76M States). Solid: steady-state graph time; dashed: achieved bandwidth.}
        \label{fig:scaling}
    \end{minipage}
\end{figure*}

\paragraph{Cross-framework timing.}
\label{sec:matrix}
\Cref{tab:matrix} and \cref{fig:hero} show that the preferred execution model
changes with game size. On the smallest trees the single-threaded C++ library wins.
The GPU pays a fixed launch-and-replay floor of about $0.113$\,ms per iteration, and
a 54-Infoset game cannot fill it. The compiled dataflow overtakes the CPU frameworks
as the game grows, and the margin widens from the HUNL turn subgame upward. There
graph replay is \liteGraphTurn$\times$ faster per iteration than LiteEFG. Across nearly four decades of game
size its per-iteration time changes by only about $5\times$. The number of dispatches
is set by tree depth, and additional states only widen each batched operation.


Kim (2026)~\cite{kim2026parallelizing} runs sequence-form CFR$^{+}$ on the same
A100, so the gap to it measures execution representation alone. The graph-replayed path is \kimGraphRange$\times$
faster across the eight-game suite, with a median of \kimGraphMedian$\times$, and the
eager compiled path is \kimEagerRange$\times$ faster. The flat schedule therefore carries most of the margin and graph replay adds the
rest. \Cref{tab:ttt} confirms the same ordering at matched exploitability
thresholds. The representation carries most of the
speedup. On eight CPU threads and no accelerator, the compiled
solver is \kimCpuRange$\times$ faster than the A100 baseline on all eight games, with
a median of \kimCpuMedian$\times$, and graph replay is a further
\graphCpuLargestRange$\times$ ahead of that CPU arm on the four largest games. Eight threads is a fair CPU ceiling for this dataflow. The sweep to 28 pinned cores
in \cref{fig:cpuscaling} saturates early, because a depth block is too small for the
intra-operator thread pool to split further. Against LiteEFG \cite{liu2024liteefg},
the general compiled CPU library, the same pattern holds. The eight-thread arm is already \cpuLiteLargestRange$\times$ faster on the four
largest games, and graph replay extends the margin to \liteGraphLargestRange$\times$.

\paragraph{Where the speed comes from.}
\label{sec:ladder}
\input{figures/TABLE_T6_overhead.tex}
The HUNL turn ablation in \cref{fig:ladder} times the two stages of the design in one
process under one software stack. Compiling the reference iteration into the
depth-level dataflow cuts Aten operations by $18\times$ and time by $19\times$. Graph
replay then removes the remaining host submissions for a further $1.6\times$, giving a
$29\times$ ladder end to end. The Nsight traces in \cref{app:profile} show the
mechanism. The eager path issues 87 kernel launches per iteration and keeps the GPU
busy a third of the time. Graph replay issues one graph launch per iteration and keeps
the GPU busy 94\% of the time. Recording 100 iterations per graph or handing the same body to
\texttt{torch.compile} \cite{paszke2019pytorch} changes nothing further. The
remaining time is the execution of the same 87 small gather and scatter kernels. The prior GPU baseline \cite{kim2026parallelizing} cannot take the second step at
all. Its iteration calls cuSPARSE through CuPy, which refuses to run under stream
capture. Its per-level row slicing and index assignment also move data between host
and device on every iteration. Capture therefore fails on the first sparse product
even at Kuhn size (\cref{app:profile}). Capturability is a property of the representation.

\Cref{fig:graphablation} shows graph-vs-eager steady-state times for every game. The
multiplier shrinks as trees grow, from 3.6$\times$ on Kuhn (54 Infosets,
fully launch-bound) to 1.6$\times$ on the HUNL turn subgame (83k Infosets), because
kernel execution time grows with data volume while launch overhead is fixed.

\begin{figure}[!htb]
  \begin{minipage}[b]{0.48\linewidth}\centering
    \includegraphics[width=\linewidth]{figures/fig7_cpu_scaling.pdf}
    \Description{Log-log plot of milliseconds per iteration against CPU thread count
    for four games; each measured curve flattens after eight threads while a dash-dot
    ideal-scaling line keeps falling.}
    \caption{CPU thread scaling of the compiled solver (eager, CFR$^{+}$); dash-dot
    lines mark ideal scaling from the one-thread point, the dotted line the 8-thread
    arm of \cref{tab:matrix}.}
    \label{fig:cpuscaling}
  \end{minipage}\hfill
  \begin{minipage}[b]{0.48\linewidth}\centering
    \includegraphics[width=\linewidth]{figures/fig3_graph_ablation.pdf}
    \Description{Grouped bar chart with one pair of bars per game on a log scale:
    a light bar for eager execution and a dark bar for CUDA-graph replay. The ratio
    printed above each pair shrinks from 3.6 times on Kuhn to 1.6 times on the HUNL
    turn subgame.}
    \caption{CUDA-graph replay vs.\ eager execution: steady-state ms per CFR$^{+}$
    iteration for every game (log scale). Labels give the eager/graph ratio.}
    \label{fig:graphablation}
  \end{minipage}
\end{figure}

The scaling sweep in \cref{fig:scaling} explains why the compiled schedule keeps
paying off on larger trees. Growing the turn subgame from 4 to 24 hands per player
multiplies the node count by $37.9\times$ while iteration time grows only
$2.95\times$. The extra hands add data volume and leave tree depth nearly unchanged,
so the operation sequence stays fixed and each operation simply gets wider. Achieved
bandwidth from a bytes-moved model reaches $5.5\%$ of the A100 peak at the largest
size, and hardware counters put measured traffic at $17\%$ (\cref{app:profile}). The
sweep is governed by launch latency and depth, not by bandwidth.

Precision and memory follow the same logic. Float64 costs $1.08\times$ on the graph
path for the turn subgame and nothing on Leduc, with eager times unchanged. Iteration
time is set by dispatch, and the arithmetic is a small part of it.
\Cref{tab:memory} (\cref{app:extra}) lists peak GPU memory per game against the
analytic model of \cref{sec:compiled}. Across the scaling sweep, allocated memory per tree node stays within
a $3\%$ band while the tree grows $37.9\times$, so memory scales linearly over the
tested range.


\begin{figure}[!htb]
    \centering
    \includegraphics[width=0.95\linewidth]{figures/fig4_curves_wide.pdf}
    \Description{Six log-log panels, one per game, plotting exploitability against
    training wall-clock seconds for up to five backends. In every panel the GPU-CFR
    CUDA-graph curve sits lowest, reaching lower exploitability earlier; the CPU arm
    of the same compiled solver follows, then the Kim baseline, LiteEFG, and
    OpenSpiel.}
    \caption{Exploitability versus training wall-clock (both axes logarithmic) on four public
    games and the native HUNL subgames.}
    \label{fig:curves}
\end{figure}

\paragraph{Quality against wall-clock.}
\label{sec:quality}
\Cref{fig:curves} plots exploitability against training wall-clock. Each backend
runs its library's native preset (\cref{tab:matrix,tab:semantics}; the update-rule
study is in \cref{sec:convergence}), so the LiteEFG and OpenSpiel traces also differ
in algorithm. The matched-semantics comparison is GPU-CFR against Kim. The HUNL
panels use the shared exports and compare GPU-CFR on CUDA, our compiled CFR$^{+}$
implementation on CPU, Kim's NoRegret CFR$^{+}$ on CUDA, and LiteEFG's native vanilla
CFR. Each arm runs 10 independent processes at six log-spaced wall-clock targets from
0.1 to 30 seconds. Points are medians and bands are the full range across processes,
which \cref{app:correctness} traces to reduction order. In every panel the
graph-replayed curve sits lowest, and the lower per-iteration cost translates
directly into time. \Cref{tab:ttt} shows GPU-CFR reaching thresholds of 10, 3 and 1.5
times Kim's 30\,s exploitability \tttRatioMin--\tttRatioMax$\times$ sooner, with
disjoint bootstrap intervals. Kuhn and Dark Hex (converged within milliseconds) and
LiteEFG's Battleship trace (accessor stuck near $10^{-15}$) are omitted.

\Cref{tab:ttt} reports, for the matched GPU-CFR versus Kim pair of \cref{fig:curves},
the distribution over independent processes of the estimated wall-clock at which each
process first crosses a fixed exploitability. Thresholds are 10, 3 and 1.5 times Kim's
30\,s median exploitability; times are read off each trace by log-log interpolation between the
bracketing checkpoints, and the intervals are 10{,}000-resample bootstraps over ten
processes. The CPU column reports the median crossing time over ten independent
processes of the same compiled solver on eight threads.
\input{figures/TABLE_T13_ttt.tex}

\paragraph{Update rule and update order.}
\label{sec:rules}
The rules differ only in how the two accumulators evolve. CFR$^{+}$
\cite{tammelin2014solving} clamps accumulated regrets at zero after each update and
weights the strategy average linearly by iteration index. DCFR
\cite{brown2019solving} instead discounts the accumulators: positive and negative
regrets decay at separate iteration-indexed rates, and the strategy average decays
polynomially. Predictive CFR$^{+}$ (PCFR$^{+}$) \cite{farina2021faster} keeps the
CFR$^{+}$ regret update but matches the strategy against the accumulated regret plus
the current instantaneous regret as a predictor. Under \emph{simultaneous} updates
both players' regrets are computed from one strategy profile; under
\emph{alternating} updates the first player's update is visible to the second
player's update in the same sweep \cite{burch2019revisiting}.

The compiled engine supports these four rules under both update orders, allowing one
execution path to expose their convergence behavior. The full grid has 96 cells. At
equal reported iteration counts, alternating updates are tighter in 44 of the 48
game/rule comparisons; the four exceptions are on Dark Hex, where both orders are
within $1.3\times10^{-6}$ of zero. After charging alternating for its approximately
two player updates, it remains better in 41 of 48 cells. The timing matrix therefore fixes
one common update order, while this study measures the algorithmic effect of
changing it.

Update order also changes which rule leads. PCFR$^{+}$ is tightest on eleven of the
twelve games under simultaneous updates, with DCFR leading only on Dark Hex. Under
alternating updates, PCFR$^{+}$ leads six games and DCFR the other six. The interaction shows that the compiler
supports a family of CFR updates without making the systems result depend on one
preferred rule. \Cref{sec:convergence} checks the ranking against LiteEFG and OpenSpiel.

\input{figures/TABLE_T9_convergence.tex}

\paragraph{Repeated solves of one tree.}
\label{sec:hunl}
\Cref{tab:lifecycle,tab:overhead} break the HUNL turn tree into setup and per-solve
phases. Building and compiling the specification dominates, graph capture is an order
of magnitude cheaper, and each later 1{,}000-iteration solve is subsecond. Setup is
repaid within the first solve. The \turnFirstSolveSec\,s from specification to
solution is below the \turnKimThousandSec\,s of Kim's A100 implementation and the
\turnLiteThousandSec\,s of LiteEFG for the same 1{,}000 iterations
(\cref{tab:matrix}).

The export mechanism extends the comparison to Libratus endgame~4
\cite{brown2018superhuman,steinberger2019pokerrl} at full range. This game has
96{,}159 Infosets on a tree of 36 million states that is almost entirely terminal
(\cref{tab:sg4threeway}). At matched iteration counts GPU-CFR is $870\times$ faster
than LiteEFG, which inverts the small-scale C++ advantage. Kim's sequence-form path
is slightly faster here, because a terminal-heavy tree leaves little depth-dependent
work for the compiled schedule to remove. The three backends agree on the value to $2.5$ chips
(\cref{app:libratus}).
\input{figures/TABLE_T8_sg4threeway.tex}
